\paragraph{Related Work}
Performance analysis strategies to provide convergence guarantees for first-order methods can be broadly split into two categories: classical and computer-assisted. Classical techniques derive analytical convergence bounds through mathematical proofs, often relying on problem-specific assumptions like smoothness or convexity to establish worst-case rates. Computer-assisted techniques, by contrast, leverage computational tools like semidefinite programming and performance estimation problems (PEPs) to numerically compute tight worst-case bounds, enabling automated verification and optimisation of algorithm performance without exhaustive manual analysis. In general, these study some performance metric $\phi^k(x)$, such as suboptimality $f(z^k(x),x) - f(z^\star(x),x)$, distance to optimality $\|z^k(x) - z^\star(x)\|$, or distance travelled $\|z^{k+1}(x) - z^k(x)\|$, for some choice of norm $\|\cdot\|$.
Guarantees are then framed either in a probabilistic sense, i.e., bounding:

\begin{equation}\label{eq:probabilistic_guarantee}
\mathbb{P}\{\phi^k(x) > \varepsilon\}
\end{equation}

for some probability measure $\mathbb{P}$ and tolerance $\varepsilon$, or in a worst-case sense, i.e., evaluating and bounding the worst-case performance $\phi^k_{\max}(x) \defeq \max_{x\in\mathcal{X}} \phi^k(x)$. Recent work has given particular attention to worst-case analysis: for instance, in \cite{VRBS25}, the authors investigated finite-step convergence verification of first-order methods for parametric convex quadratic programs, obtaining tight, data-free quantifications of $\phi^k_{\max}(x)$ that incorporate warm-starting effects; similarly, \cite{VRBS+24} formulated an exact mixed-integer linear program to maximise $\phi^k(x)$, defined using the $\infty$-norm, by encoding the algorithm’s dynamics as polyhedral constraints. While these approaches yield rigorous guarantees, they tend to be conservative for typical parameter instances. To address these limitations, probabilistic analyses have emerged as a promising alternative: for instance, \cite{RSBS25} developed probabilistic guarantees for both classical and learned optimisers using sample convergence bounds and the Probably Approximately Correct (PAC)-Bayes framework, respectively, while \cite{JHKM+25} employed advancements in the so-called Scenario Approach Theory to generate tight, data-driven guarantees for both convergence and general definitions of $\phi^k(x)$. However, these approaches suffer from important drawbacks: the PAC-Bayes framework is unintuitive and potentially limited in proving generalisation for certain learnable tasks, whereas Scenario Approach bounds are heavily sample-count dependent, which can lead to overly conservative guarantees with insufficient data.